\chapter{Workbook Phần 1: 20 Bài Tập Thực Hành Excel \& Báo Cáo Nghiệp Vụ}

\begin{lythuyetbox}[Mục Tiêu Luyện Tập Phần 1]
Phần này cung cấp \textbf{20 bài tập thực hành Excel chuyên sâu} từ cơ bản đến nâng cao. Mỗi bài tập đều gắn liền với tệp dữ liệu cụ thể, mô phỏng tình huống thực tế tại các phòng ban (Tài chính, Kinh doanh, Marketing, Vận hành), đi kèm công thức chuẩn và kết quả số liệu chính xác để bạn tự kiểm tra.
\end{lythuyetbox}

\begin{thuchanhbox}[Tệp Dữ Liệu Thực Hành Cụ Thể (Bấm Vào Đường Dẫn Để Mở File)]
Toàn bộ 20 bài tập trong chương này được liên kết trực tiếp với bảng tính thực tế đã chuẩn bị sẵn trong kho dữ liệu:
\begin{center}
\href{run:./data/excel_practice_data.xlsx}{\Large\textbf{\color{navyprimary}data/excel\_practice\_data.xlsx}} \\[0.2cm]
\small\textit{(Nhấp chuột vào liên kết trên để mở file Excel trên máy tính của bạn, hoặc xem trực tuyến tại: \url{https://github.com/TontonYuta/DataLearning/blob/main/data/excel_practice_data.xlsx})}
\end{center}
Tệp gồm 6 Sheets chuyên biệt: \texttt{NhanVien} (100 dòng), \texttt{CuocVanChuyen} (Ma trận 8 tỉnh), \texttt{DoanhSo} (5,000 giao dịch), \texttt{ChuoiHoTen} (50 mẫu tên), \texttt{TidyData\_Unpivot} (20 sản phẩm), \texttt{XuatKho} (200 lượt xuất hàng).
\end{thuchanhbox}

\begin{baitapbox}[Hệ Thống 20 Bài Tập Excel Thực Chiến Với Số Liệu Cụ Thể]
\begin{enumerate}
    \item \textbf{Bài 1 (XLOOKUP Cơ Bản -- Sheet \texttt{NhanVien})}: Mở tệp \href{run:./data/excel_practice_data.xlsx}{\texttt{data/excel\_practice\_data.xlsx}}, chọn sheet \texttt{NhanVien}. Bảng gồm các cột \texttt{A:D} là \texttt{[Ma\_NV, Ho\_Ten, Phong\_Ban, Luong\_VND]}. Tại ô \texttt{F2}, nhập mã nhân viên \texttt{"NV105"}. Viết công thức XLOOKUP tại ô \texttt{G2} để tra cứu Họ Tên và tại \texttt{H2} để tra cứu Lương, trả về chuỗi \texttt{"Không tìm thấy"} nếu mã không tồn tại. \textit{Kết quả kiểm tra}: Họ tên là \textbf{"Phan Van Dat"}, Lương là \textbf{28,000,000 VND}.
    \item \textbf{Bài 2 (XLOOKUP Hai Chiều -- Sheet \texttt{CuocVanChuyen})}: Mở sheet \texttt{CuocVanChuyen} chứa ma trận cước phí gồm cột A là Tỉnh/Thành đến và các cột \texttt{B:E} là mức cân nặng (\texttt{1kg, 2kg, 5kg, 10kg}). Nhập Tỉnh đến tại \texttt{G2 = "Da Nang"} và Khối lượng tại \texttt{H2 = "5kg"}. Viết công thức 2-Way XLOOKUP tra cứu chính xác cước phí tại \texttt{I2}. \textit{Kết quả kiểm tra}: Cước phí là \textbf{55,000 VND}.
    \item \textbf{Bài 3 (INDEX \& MATCH Hai Chiều -- Sheet \texttt{NhanVien})}: Tại sao trong hệ thống lớn hoặc phiên bản cũ, cặp hàm \texttt{INDEX/MATCH} lại an toàn hơn \texttt{VLOOKUP}? Viết công thức \texttt{INDEX/MATCH} tra cứu Họ Tên ở cột B dựa trên Mã NV ở ô \texttt{F2} mà không làm lệch cột khi có người chèn thêm cột mới vào bảng.
    \item \textbf{Bài 4 (SUMIFS Đa Điều Kiện -- Sheet \texttt{DoanhSo})}: Mở sheet \texttt{DoanhSo} gồm 5,000 dòng giao dịch (\texttt{A:G}). Viết công thức \texttt{SUMIFS} tính tổng doanh thu thỏa mãn đồng thời 3 điều kiện: Khu vực (\texttt{Khu\_Vuc}) là \texttt{"Mien Bac"}, Ngành hàng (\texttt{Nganh\_Hang}) là \texttt{"Dien Tu"}, và phát sinh trong Quý 1 năm 2026 (từ ngày \texttt{>=2026-01-01} đến \texttt{<=2026-03-31}).
    \item \textbf{Bài 5 (COUNTIFS Đếm Đơn Hàng -- Sheet \texttt{DoanhSo})}: Viết công thức \texttt{COUNTIFS} đếm số lượng đơn hàng có giá trị \texttt{Doanh\_Thu\_USD > 500} của nhân viên \texttt{"Nguyen Van A"} có trạng thái đơn hàng là \texttt{"COMPLETED"}.
    \item \textbf{Bài 6 (Dynamic Array UNIQUE \& SORT -- Sheet \texttt{DoanhSo})}: Tại sheet \texttt{DoanhSo}, cột C chứa danh sách nhân viên bị lặp lại nhiều lần. Viết công thức mảng động xuất ra danh sách nhân viên duy nhất và tự động sắp xếp theo thứ tự bảng chữ cái A--Z. \textit{Kết quả kiểm tra}: Trả về mảng động đúng 6 nhân viên.
    \item \textbf{Bài 7 (Dynamic Array FILTER -- Sheet \texttt{DoanhSo})}: Viết công thức mảng động \texttt{FILTER} trích xuất toàn bộ các dòng từ bảng doanh số thỏa mãn 2 điều kiện: \texttt{Trang\_Thai = "COMPLETED"} và \texttt{Doanh\_Thu\_USD >= 2000}.
    \item \textbf{Bài 8 (Tính Hoa Hồng Đa Bậc IFS)}: Chính sách hoa hồng doanh số tháng: Dưới hoặc bằng \$1,000: hoa hồng 3\%; Từ \$1,001 đến \$2,000: hoa hồng 5\%; Trên \$2,000: hoa hồng 8\%. Viết công thức \texttt{IFS} tính tiền hoa hồng cho đơn hàng tại ô \texttt{F2} có giá trị \$1,500. \textit{Kết quả kiểm tra}: Hoa hồng là \textbf{\$75.00}.
    \item \textbf{Bài 9 (Tách Chuỗi Họ Tên -- Sheet \texttt{ChuoiHoTen})}: Mở sheet \texttt{ChuoiHoTen} có 50 dòng họ và tên đầy đủ ở cột A. Sử dụng các hàm xử lý chuỗi hiện đại \texttt{TEXTBEFORE} và \texttt{TEXTAFTER} để tách riêng Tên gọi sang cột B và Họ đệm sang cột C.
    \item \textbf{Bài 10 (Chuẩn Hóa Ngày Tháng Bằng Power Query)}: Cột A chứa ngày tháng ở các định dạng hỗn loạn: \texttt{"20260315"}, \texttt{"15/03/2026"}, \texttt{"2026.03.15"}. Trình bày quy trình dùng Power Query qua chức năng \textit{Using Locale} để chuẩn hóa toàn bộ về định dạng chuẩn \texttt{Date (DD/MM/YYYY)}.
    \item \textbf{Bài 11 (Tính Tăng Trưởng MoM \& Bẫy Lỗi IFERROR)}: Cho cột doanh thu 12 tháng tại \texttt{B2:B13}. Viết công thức tính tỷ lệ tăng trưởng phần trăm so với tháng trước (\% MoM) tại ô \texttt{C3} và dùng \texttt{IFERROR} bẫy lỗi chia cho 0 nếu tháng trước doanh thu bằng 0.
    \item \textbf{Bài 12 (Tự Động Đánh Số Thứ Tự Động)}: Viết công thức đánh số thứ tự 1, 2, 3... tại cột A sao cho số thứ tự chỉ xuất hiện khi cột B có nhập dữ liệu, và tự động bỏ trống nếu cột B rỗng.
    \item \textbf{Bài 13 (Unpivot Columns Trong Power Query -- Sheet \texttt{TidyData\_Unpivot})}: Mở sheet \texttt{TidyData\_Unpivot} gồm cột \texttt{Ma\_SP} và 12 cột tháng (\texttt{Thang\_1} đến \texttt{Thang\_12}). Trình bày các bước Unpivot Columns trong Power Query để đưa dữ liệu ngang về chuẩn Tidy Data gồm 3 cột: \texttt{[Ma\_SP, Thang, Doanh\_Thu]}.
    \item \textbf{Bài 14 (Pivot Table Lọc Top 5 Nhân Viên Doanh Số Cao Nhất)}: Dựa trên dữ liệu sheet \texttt{DoanhSo}, trình bày các bước tạo Pivot Table và thiết lập bộ lọc tự động \textit{Value Filters $\rightarrow$ Top 5} để hiển thị 5 nhân viên có tổng doanh thu lớn nhất.
    \item \textbf{Bài 15 (Pivot Table Calculated Field)}: Trong Pivot Table, tạo một trường tính toán mới mang tên \texttt{"Loi\_Nhuan\_Uoc\_Tinh"} với công thức: \texttt{'Doanh\_Thu\_USD' * 0.30 - 15}.
    \item \textbf{Bài 16 (Xác Thực Dữ Liệu Data Validation)}: Thiết lập ô \texttt{D2} chỉ cho phép người dùng nhập ngày trong tương lai (sau ngày hiện tại \texttt{TODAY()}) và hiển thị thông báo lỗi tùy chỉnh: \textit{"Ngày nhập vào phải sau ngày hôm nay!"}.
    \item \textbf{Bài 17 (Conditional Formatting Tô Màu Cả Dòng)}: Tại sheet \texttt{DoanhSo}, thiết lập định dạng có điều kiện tự động tô màu nền đỏ nhạt cho toàn bộ dòng dữ liệu từ cột A đến G nếu cột Trạng thái (\texttt{Trang\_Thai}) là \texttt{"OVERDUE"}.
    \item \textbf{Bài 18 (XLOOKUP Tìm Kiếm Ngược Từ Dưới Lên -- Sheet \texttt{XuatKho})}: Mở sheet \texttt{XuatKho} có 200 dòng lịch sử xuất hàng. Viết công thức XLOOKUP với tham số \texttt{search\_mode = -1} để tìm Đơn giá xuất kho của \textbf{lần xuất gần đây nhất} của sản phẩm \texttt{"SP\_IPHONE\_15"}.
    \item \textbf{Bài 19 (Tính Số Ngày Làm Việc Thực Tế NETWORKDAYS.INTL)}: Viết công thức tính số ngày làm việc thực tế giữa ngày bắt đầu tại \texttt{A2 = "2026-01-01"} và ngày kết thúc tại \texttt{B2 = "2026-01-31"}, tự động loại trừ Thứ Bảy, Chủ Nhật và danh sách các ngày nghỉ Tết tại dải ô \texttt{Holiday\_List}.
    \item \textbf{Bài 20 (Ghép Nối Chuỗi Hàng Loạt TEXTJOIN)}: Cho cột danh sách email nhân viên từ \texttt{C2:C20}. Viết công thức ghép toàn bộ các email này thành một chuỗi duy nhất phân tách bằng dấu chấm phẩy và khoảng trắng (\texttt{; }) để dán trực tiếp vào mục gửi thư Outlook, tự động bỏ qua các ô trống.
\end{enumerate}
\end{baitapbox}

\begin{dapanbox}[Đáp Án \& Lời Giải Ngắn Gọn Cho 20 Bài Tập Excel]
\begin{itemize}
    \item \textbf{Bài 1}: Tại \texttt{G2 = XLOOKUP(F2, A2:A101, B2:B101, "Không tìm thấy")}; Tại \texttt{H2 = XLOOKUP(F2, A2:A101, D2:D101, "Không tìm thấy")}. Kết quả: \textbf{"Phan Van Dat"} và \textbf{28000000}.
    \item \textbf{Bài 2}: \texttt{=XLOOKUP(G2, A2:A9, XLOOKUP(H2, B1:E1, B2:E9))}. Kết quả: \textbf{55,000 VND}.
    \item \textbf{Bài 3}: \texttt{=INDEX(B2:B101, MATCH(F2, A2:A101, 0))}. Ưu điểm: Không bị hardcode số thứ tự cột như VLOOKUP (vốn sẽ tính sai nếu người dùng chèn thêm cột mới).
    \item \textbf{Bài 4}: \texttt{=SUMIFS(F2:F5001, D2:D5001, "Mien Bac", E2:E5001, "Dien Tu", B2:B5001, ">=2026-01-01", B2:B5001, "<=2026-03-31")}.
    \item \textbf{Bài 5}: \texttt{=COUNTIFS(C2:C5001, "Nguyen Van A", F2:F5001, ">500", G2:G5001, "COMPLETED")}.
    \item \textbf{Bài 6}: \texttt{=SORT(UNIQUE(C2:C5001))}. Trả về mảng động 6 nhân viên theo thứ tự A--Z.
    \item \textbf{Bài 7}: \texttt{=FILTER(A2:G5001, (G2:G5001="COMPLETED") * (F2:F5001>=2000), "Không có đơn hàng")}.
    \item \textbf{Bài 8}: \texttt{=IFS(F2<=1000, F2*0.03, F2<=2000, F2*0.05, F2>2000, F2*0.08)}. Kết quả: \textbf{75.00 USD}.
    \item \textbf{Bài 9}: Tên gọi: \texttt{=TEXTAFTER(A2, " ", -1)}. Họ đệm: \texttt{=TEXTBEFORE(A2, " ", -1)}.
    \item \textbf{Bài 10}: Trong Power Query: Chuột phải vào cột ngày $\rightarrow$ \textit{Change Type} $\rightarrow$ \textit{Using Locale...} $\rightarrow$ Data Type: \textit{Date}, Locale: chọn định dạng tương ứng của file nguồn.
    \item \textbf{Bài 11}: \texttt{=IFERROR((B3 - B2) / B2, 0)}. Định dạng ô hiển thị dạng Percentage (\%).
    \item \textbf{Bài 12}: \texttt{=IF(B2<>"", COUNTA(\$B\$2:B2), "")}.
    \item \textbf{Bài 13}: Mở Power Query $\rightarrow$ Chọn cột \texttt{Ma\_SP} $\rightarrow$ tab \textit{Transform} $\rightarrow$ bấm mũi tên cạnh \textit{Unpivot Columns} $\rightarrow$ chọn \textit{Unpivot Other Columns}.
    \item \textbf{Bài 14}: Kéo \texttt{Nhan\_Vien} vào Rows, \texttt{Doanh\_Thu\_USD} vào Values $\rightarrow$ Chuột phải vào 1 tên nhân viên $\rightarrow$ \textit{Filter} $\rightarrow$ \textit{Top 10...} $\rightarrow$ chọn Top 5 by Sum of Doanh\_Thu\_USD.
    \item \textbf{Bài 15}: Chọn PivotTable $\rightarrow$ Tab \textit{PivotTable Analyze} $\rightarrow$ \textit{Fields, Items, \& Sets} $\rightarrow$ \textit{Calculated Field} $\rightarrow$ Name: \texttt{Loi\_Nhuan\_Uoc\_Tinh}, Formula: \texttt{=Doanh\_Thu\_USD * 0.3 - 15}.
    \item \textbf{Bài 16}: Chọn \texttt{D2} $\rightarrow$ Tab \textit{Data} $\rightarrow$ \textit{Data Validation} $\rightarrow$ Allow: \textit{Date} $\rightarrow$ Data: \textit{greater than} $\rightarrow$ Start date: \texttt{=TODAY()}. Tab \textit{Error Alert}: nhập tiêu đề và thông báo lỗi.
    \item \textbf{Bài 17}: Chọn toàn bộ bảng \texttt{\$A\$2:\$G\$5001} $\rightarrow$ \textit{Conditional Formatting} $\rightarrow$ \textit{New Rule} $\rightarrow$ \textit{Use a formula...} $\rightarrow$ Nhập: \texttt{=\$G2="OVERDUE"} (khóa cột G bằng dấu \$).
    \item \textbf{Bài 18}: \texttt{=XLOOKUP("SP\_IPHONE\_15", B2:B201, D2:D201, "N/A", 0, -1)}. Tham số \texttt{-1} ở vị trí thứ 6 chỉ định tìm từ dưới lên trên.
    \item \textbf{Bài 19}: \texttt{=NETWORKDAYS.INTL(A2, B2, 1, Holiday\_List)} (tham số 1 tương ứng nghỉ Thứ 7 và Chủ Nhật).
    \item \textbf{Bài 20}: \texttt{=TEXTJOIN("; ", TRUE, C2:C20)} (tham số \texttt{TRUE} tự động bỏ qua các ô rỗng).
\end{itemize}
\end{dapanbox}
